\documentclass[11pt]{article}
\usepackage[margin=1in]{geometry}
\usepackage{amsmath,amsthm}
\usepackage{stix2}
\usepackage{microtype}
\usepackage{enumitem}
\usepackage[hidelinks]{hyperref}

\DeclareMathOperator{\Hom}{Hom}
\DeclareMathOperator{\im}{im}
\DeclareMathOperator{\coim}{coim}
\DeclareMathOperator{\coker}{coker}
\DeclareMathOperator{\Ch}{Ch}
\newcommand{\Z}{\mathbb{Z}}
\newcommand{\Q}{\mathbb{Q}}
\newcommand{\R}{\mathbb{R}}
\newcommand{\id}{\mathrm{id}}

\setlist[enumerate]{label=(\alph*),ref=\alph*,leftmargin=*,itemsep=.45em}
\renewcommand{\thesection}{Exercise \arabic{section}.}
\setlength{\parindent}{0pt}
\setlength{\parskip}{.25em}
\hypersetup{pdftitle={Algebraic Topology I: Homework 1}}

\begin{document}
\begin{center}
  {\Large\bfseries Algebraic Topology I: Homework 1}\\[.4em]
  Additive and abelian categories, chain complexes, and chain homotopies
\end{center}

Throughout, an \emph{abelian category} means an additive category with all
kernels and cokernels in which every monomorphism is a kernel and every
epimorphism is a cokernel. In an arbitrary abelian category, give arguments
using morphisms and universal properties.

Use homological grading: a chain complex has differentials
\(d_n:C_n\to C_{n-1}\) satisfying \(d_{n-1}d_n=0\).
A chain homotopy between chain maps \(f,g:C\to D\) consists of morphisms
\(h_n:C_n\to D_{n+1}\) satisfying
\[
f_n-g_n=d^D_{n+1}h_n+h_{n-1}d^C_n.
\]

\section{Mono, epi, and iso}
Let \(\mathcal C\) be a category. A morphism \(f:X\to Y\) is called:
\begin{itemize}[leftmargin=*]
\item a \emph{monomorphism} if, for every object \(T\) and morphisms
\(u,v:T\to X\), the equality \(fu=fv\) implies \(u=v\);
\item an \emph{epimorphism} if, for every object \(T\) and morphisms
\(u,v:Y\to T\), the equality \(uf=vf\) implies \(u=v\);
\item an \emph{isomorphism} if there exists a morphism \(g:Y\to X\)
such that \(gf=1_X\) and \(fg=1_Y\).
\end{itemize}

\begin{enumerate}
\item Prove that the inverse of an isomorphism is unique, and that every
isomorphism is both a monomorphism and an epimorphism. In the category of
sets, prove that monomorphisms are precisely injections, epimorphisms are
precisely surjections, and isomorphisms are precisely bijections.

\item Let \(\mathbf{CRing}\) be the category of commutative rings with
identity and identity-preserving ring homomorphisms. Show that the inclusion
\[
\Z\hookrightarrow\Q
\]
is both a monomorphism and an epimorphism, but is not an isomorphism. Thus
an epimorphism of rings need not be surjective.

\item Let \(\mathbf{Haus}\) be the category of Hausdorff topological spaces
and continuous maps. Give \(\Q\) and \(\R\) their usual topologies. Show
that the inclusion
\[
\Q\hookrightarrow\R
\]
is both a monomorphism and an epimorphism in \(\mathbf{Haus}\), but is not
an isomorphism. Is this inclusion still an epimorphism in the category
\(\mathbf{Top}\) of all topological spaces and continuous maps? Justify
your answer.

\item Let \(\mathcal A\) be an abelian category. Prove that a morphism
in \(\mathcal A\) is an isomorphism if and only if it is both a
monomorphism and an epimorphism.
\end{enumerate}

\section{Coimage and image}
Let \(f:X\to Y\) be a morphism in an abelian category. Define
\[
\coim(f)=\coker(\ker f),
\qquad
\im(f)=\ker(\coker f).
\]
Let \(p:X\to\coim(f)\) and \(i:\im(f)\to Y\) be the canonical morphisms.
\begin{enumerate}
\item Prove that there is a unique morphism
\[
\bar{f}:\coim(f)\longrightarrow\im(f),
\qquad
f=i\circ\bar{f}\circ p.
\]
\item Prove that \(\bar{f}\) is an isomorphism.
\end{enumerate}

\section{Matrix algebra in additive categories}
Let \(\mathcal C\) be an additive category. For a biproduct
\(X_1\oplus X_2\), write
\[
\iota_j^X:X_j\longrightarrow X_1\oplus X_2,
\qquad
\rho_i^X:X_1\oplus X_2\longrightarrow X_i
\]
for its canonical inclusion and projection maps. The projections are the
product structure maps, and the inclusions are characterized by
\[
\rho_i^X\iota_j^X=
\begin{cases}
1_{X_i},&i=j,\\
0,&i\ne j.
\end{cases}
\]
The inclusions also make \(X_1\oplus X_2\) a coproduct. Use analogous
notation for \(Y_1\oplus Y_2\).

\begin{enumerate}
\item Show that a morphism
\[
f:X_1\oplus X_2\longrightarrow Y_1\oplus Y_2
\]
is uniquely determined by its four components
\[
f_{ij}=\rho_i^Y f\iota_j^X.
\]
Express \(f\) in terms of these components, and prove that composition
corresponds to matrix multiplication.

\item Let \(\mathcal F\) be the category of finitely generated free
abelian groups, with all group homomorphisms as morphisms. Show that
\(\mathcal F\) is additive. Prove that
\[
2:\Z\longrightarrow\Z
\]
is both a monomorphism and an epimorphism in \(\mathcal F\), but is not
an isomorphism.

\item Compute the kernel and cokernel of this morphism \emph{in
\(\mathcal F\)}. Then compute its image, coimage, and the canonical morphism
\[
\coim(2)\longrightarrow\im(2).
\]
Explain why \(\mathcal F\) is not abelian.
\end{enumerate}

\section{Splittings and idempotents}
Let
\[
0\longrightarrow A\xrightarrow{i}B\xrightarrow{p}C\longrightarrow0
\]
be a short exact sequence in an abelian category \(\mathcal A\).
\begin{enumerate}
\item Suppose \(s:C\to B\) satisfies \(ps=1_C\). Construct \(r:B\to A\)
such that
\[
ri=1_A,\qquad rs=0,\qquad ir+sp=1_B.
\]
Deduce that
\[
[\,i\ \ s\,]:A\oplus C\longrightarrow B
\]
is an isomorphism, and describe its inverse.

\item Prove that \(p\) has a section if and only if \(i\) has a retraction.
Thus either condition characterizes a split short exact sequence.

\item Let \(e:X\to X\) satisfy \(e^2=e\). Construct an isomorphism
\[
X\cong\im(e)\oplus\ker(e)
\]
under which \(e\) has matrix
\[
\begin{pmatrix}1&0\\0&0\end{pmatrix}.
\]
\end{enumerate}

\section{Chain complexes form an abelian category}
Let \(\mathcal A\) be an abelian category, and write
\(\Ch(\mathcal A)\) for its category of chain complexes and chain maps.
\begin{enumerate}
\item Describe addition of chain maps and finite biproducts of complexes.
Given a chain map \(f:C\to D\), construct complexes \(K\) and \(Q\) with
\[
K_n=\ker(f_n),\qquad Q_n=\coker(f_n),
\]
and prove that they are respectively the kernel and cokernel of \(f\) in
\(\Ch(\mathcal A)\).

\item Prove that \(\Ch(\mathcal A)\) is abelian. Show that a sequence
of complexes
\[
0\longrightarrow C\longrightarrow D\longrightarrow E\longrightarrow0
\]
is short exact if and only if it is short exact in every degree.

\item Put
\[
Z_n(C)=\ker(d_n),\qquad B_n(C)=\im(d_{n+1}).
\]
Construct the canonical monomorphism \(B_n(C)\to Z_n(C)\), and define
\[
H_n(C)=\coker\bigl(B_n(C)\longrightarrow Z_n(C)\bigr).
\]
Show that chain maps induce morphisms on homology and that
\[
H_n:\Ch(\mathcal A)\longrightarrow\mathcal A
\]
is an additive functor.
\end{enumerate}

\section{What chain homotopy remembers}
Let \(C,D\) be chain complexes in an abelian category \(\mathcal A\).
\begin{enumerate}
\item Prove that chain homotopy is an equivalence relation on the set of
chain maps \(C\to D\). Show that it is compatible with addition and
composition. Deduce that complexes and homotopy classes of chain maps
form an additive category \(K(\mathcal A)\).

\item Prove that homotopic chain maps induce the same morphisms on
homology. Deduce that a chain homotopy equivalence is a
\emph{quasi-isomorphism}, meaning a chain map inducing isomorphisms on
all homology objects.

\item Work now in abelian groups. Let \(C\) be
\[
0\longrightarrow\Z\xrightarrow{\,2\,}\Z\longrightarrow0,
\]
concentrated in degrees \(1,0\), and let \(D\) have \(\Z\) in degree
\(1\) and zero elsewhere.

Describe all chain maps \(C\to D\), and determine exactly when two are
chain homotopic. Compute the abelian group
\[
\Hom_{K(\mathbf{Ab})}(C,D).
\]
Show that every chain map \(C\to D\) induces zero on homology, although
some are not nullhomotopic.
\end{enumerate}

\section{Acyclic versus contractible}
A complex \(C\) is \emph{acyclic} if \(H_n(C)=0\) for every \(n\), and
\emph{contractible} if \(1_C\) is chain homotopic to zero.
\begin{enumerate}
\item Prove that \(C\) is contractible if and only if it is acyclic and
every short exact sequence
\[
0\longrightarrow Z_n(C)\longrightarrow C_n
\xrightarrow{\bar{d}_n}Z_{n-1}(C)\longrightarrow0
\]
splits. Here \(\bar{d}_n\) is the factorization of \(d_n\) through
\(Z_{n-1}(C)\); acyclicity makes it an epimorphism.

\item Consider
\[
0\longrightarrow\Z\xrightarrow{\,2\,}\Z
\xrightarrow{\mathrm{mod}\,2}\Z/2\longrightarrow0,
\]
with the displayed nonzero terms in degrees \(2,1,0\). Show that this
complex is acyclic but not contractible. Deduce that a quasi-isomorphism
need not be a chain homotopy equivalence.

\item Let \(C\) be a chain complex of vector spaces over a field. Show
that \(C\) is chain homotopy equivalent to the complex having \(H_n(C)\)
in degree \(n\) and all differentials zero. Explain where choices enter
your construction.

\item Work with modules over the ring \(\Z\). Consider the complex
\[
C:\qquad 0\longrightarrow\Z\xrightarrow{\,2\,}\Z\longrightarrow0,
\]
concentrated in degrees \(1,0\). Compute its homology, and let \(H(C)\)
denote its homology regarded as a complex with zero differentials.

Construct a quasi-isomorphism \(C\to H(C)\), but prove that \(C\) and
\(H(C)\) are not chain homotopy equivalent. Thus the conclusion of
part~(c) can fail for modules over a ring that is not a field, even when
every term of \(C\) is a free module.
\end{enumerate}

\end{document}
